\documentclass[12pt,a4paper,english]{article}

\usepackage[utf8]{inputenc}
\usepackage[a4paper,hmargin=10mm,vmargin=25mm]{geometry}
\usepackage[backend=biber,style=numeric,sortcites,firstinits=true]{biblatex}
\usepackage[skip=5pt, indent=0pt]{parskip}
\usepackage[hidelinks]{hyperref}
\usepackage[english]{babel}
\usepackage[T1]{fontenc}
\usepackage{indentfirst}
\usepackage{enumitem}
\usepackage{textcomp}
\usepackage{url}
\usepackage{csquotes}
\usepackage{graphicx}
\usepackage{float}
\usepackage{amsmath}
\usepackage{amssymb}
\usepackage{amsfonts}

\urlstyle{sf}
\addbibresource{./refs.bib}
\graphicspath{ {../figures/} }

\title{\textbf{OCES5303} \linebreak Extended Project}
\author{Jonas Mathisrud Sterud}

\begin{document}
    \maketitle

    \begin{abstract}
        Abstract here!
    \end{abstract}

    \begin{figure}[H]
        \begin{center}
            \includegraphics[width=15cm]{cover.jpg}
        \end{center}
    \end{figure}

    \newpage
    \tableofcontents
    \newpage

    \section{The Data}

    \subsection{Raw data}

    We have five different folders with data:
    \begin{description}[labelindent=1cm]
        \item[TRAIN\_png] - Dated \verb|.png| training images across various categories.
        \item[TEST\_png] - Dated \verb|.png| testing images across various categories.
        \item[LABELS\_TRAIN] - Dated \verb|.csv| files with training labels.
        \item[train\_images\_pkl] - Dated folders, each containing various \verb|.pkl| files for the traning dataset.
        \item[test\_image\_pkl] - Dated folders, each containing various \verb|.pkl| files for the test dataset.
    \end{description}
    
    The \verb|.pkl| files contain various useful information, such as the contours and centres of the cyclonic and anticyclonic eddy's (in the training set).
    They also contain latitude and longitude positions of the images, which we'll use later for visualization.

    Additionaly, we have labels for the test dataset as a \verb|.pkl| file.

    \subsection{Parsing}

    We load, and parse, all the \verb|.pkl| files and put them into a dictionary.
    The \verb|.png| images are loaded in as grayscale images, and are then sorted under their respective dates, categories and image numbers - which are extracted from their filenames.
    
    Then, we combine all the images, and their corresponding metadata (from the \verb|.pkl| files) into a Pandas \verb|Dataframe|.

    \subsection{Analysis}



%    \tableofcontents
%    \newpage



%    \newpage
%    \printbibliography
\end{document}
